\documentclass[11pt]{article}
\usepackage[T1]{fontenc}
\usepackage{lmodern,amsmath,amssymb,geometry,microtype,xcolor,fancyhdr,hyperref}
\geometry{letterpaper,left=0.95in,right=0.95in,top=0.9in,bottom=0.88in}
\definecolor{ink}{HTML}{183953}\definecolor{tone}{HTML}{386475}
\hypersetup{colorlinks=true,urlcolor=tone,pdftitle={Sharp continuity of mutual information on a fixed-marginal fibre}}
\setlength{\parindent}{0pt}\setlength{\parskip}{5pt}
\pagestyle{fancy}\fancyhf{}\setlength{\headheight}{15pt}
\lhead{\small\sffamily\color{tone} QUANTUM INFORMATION}\rhead{\small\sffamily\color{tone} DRAFT FOR REVIEW}
\rfoot{\small\sffamily\color{tone}\thepage}
\newcommand{\shead}[1]{\par\vspace{3pt}{\large\bfseries\sffamily\color{ink}#1}\par}
\newcommand{\Tr}{\operatorname{Tr}} \newcommand{\id}{\mathrm{id}} \newcommand{\1}{\mathbf 1}
\newcommand{\C}{\mathcal C}\newcommand{\D}{\mathcal D}\newcommand{\N}{\mathcal N}\newcommand{\M}{\mathcal M}
\begin{document}
{\small\sffamily\color{tone}18 NOVEMBER 2026 \quad / \quad WORKING DRAFT}\par
{\LARGE\bfseries\sffamily\color{ink} Sharp continuity of mutual information\\[2pt] when one marginal is fixed}\par
\vspace{5pt}
Consider two bipartite quantum states with a common marginal, not assumed to be the smaller subsystem. The exact dimension-only uniform modulus involves the smaller of the two dimensions. The additional step beyond the sharp conditional-entropy estimate is to choose a product reference adapted to the anchor state and cancel its local logarithm using the common marginal.
\shead{1. The statement and its scope}
Let $A,B$ be finite-dimensional complex Hilbert spaces of dimensions $a,b$; let $r=\min(a,b)$. Write $T(\rho,\sigma)=\tfrac12\|\rho-\sigma\|_1$,
$H(\rho)=-\Tr\rho\log\rho$, $H(A|B)=H(AB)-H(B)$, and
$I(A:B)=H(A)+H(B)-H(AB)$. All logarithms are natural. Set $h(t)=-t\log t-(1-t)\log(1-t)$ with $0\log0=0$. For $K>1$ let
\[
g_K(\delta)=\begin{cases}
 h(\delta)+\delta\log(K-1),&0\le\delta\le1-K^{-1},\\
 \log K,&1-K^{-1}\le\delta\le1.
\end{cases}
\qquad g_1\equiv0 .
\]
\textbf{Theorem.} In every fixed pair of dimensions and for every $0\le\delta\le1$,
\begin{equation*}
 \max_{\substack{\rho_{AB},\sigma_{AB}\text{ states},\ \rho_A=\sigma_A\\T(\rho,\sigma)\le\delta}}
 |I(A:B)_\rho-I(A:B)_\sigma|=g_{r^2}(\delta).
 \tag{1}
\end{equation*}
The same holds on exchanging $A$ and $B$. The common marginal varies over the optimization: (1) does \emph{not} promise sharpness for each preassigned marginal. Both marginals may coincide. For $r\ge2$, $\delta>0$, even the subclass in which the other marginal must differ has the same \emph{supremum}. No maximizer within that subclass is asserted.
\shead{2. What was already known}
Berta, Lami and Tomamichel [1, \S IV] ask about mutual information with one common marginal just before posing the stronger arbitrary-pair question (106). The classical counterexample in [2] changes both marginals; (1) does not restore (106) without a marginal constraint.

A weaker estimate already used the smaller dimension even when the fixed marginal was on the larger side. Shirokov [6, Cor.\ 1, after (25)], with trivial conditioning system, obtained $2\delta\log r+(1+\delta)h(\delta/(1+\delta))$. The same bound appears for entanglement-assisted classical capacity in [6, Prop.\ 10, (81)]. That work already uses the isotropic comparison and the identity/depolarizing channel pair.

Berta et al. [3, Thm.\ 1.1] give the sharp conditional-entropy estimate in terms of the dimension $a$ of $A$ in $H(A|B)$. This settles (1) if $a\le b$, since $H(A)$ cancels. Their Prop.\ 2.1 gives a dimension-only bound $g_{ar}$ without a fixed marginal. Their Prop.\ 3.1 and Rem.\ 3.2 settle both the capacity bound and sharpness for input dimension at most output dimension. The case addressed here is $a>b$: the exact curve is $g_{b^2}$, for every fixed pair of dimensions and every distance budget. We combine the known max-information dimension bound [4, Lem.\ B.7] with the interpolation of [3].

\textbf{Status.} OpenAI AI assistance was used for the argument, drafting and internal checks. This is a working proof, not an external peer review or a formalization. No historical priority is claimed.
\newpage
\shead{3. A product reference adapted to the anchor}
We first work with positive definite $\rho,\sigma$. For every $Y\ge0$ on $A\otimes B$,
\[
 Y\ \le\ b(\Tr_B Y)\otimes \1_B,\qquad
 Y\ \le\ a\,\1_A\otimes(\Tr_A Y).
 \tag{2}
\]
For the first inequality, write a rank-one summand as
$|y\rangle\langle y|$, with $y=\sum_{j=1}^{b}y_j\otimes e_j$. For $v=\sum_j v_j\otimes e_j$,
\[
 |\langle y,v\rangle|^2\le b\sum_j|\langle y_j,v_j\rangle|^2
 \le b\sum_j\langle v_j,(\textstyle\sum_k |y_k\rangle\langle y_k|)v_j\rangle .
\]
The last matrix in parentheses is $\Tr_B |y\rangle\langle y|$. Sum rank-one summands for general $Y$; interchange the systems for the other bound.

Set $K=r^2$. There exists a positive definite density operator $\eta_A$ such that
\begin{equation*}
       \sigma\le K W,\qquad W=\eta_A\otimes\sigma_B.\tag{3}
\end{equation*}
Indeed, if $a\le b$, choose $\eta_A=\1_A/a$ and apply (2). If $b<a$, set
\[
  X=(\1_A\otimes\sigma_B^{-1/2})\sigma(\1_A\otimes\sigma_B^{-1/2}),
  \qquad \eta_A=\frac1b\Tr_B X .
\]
Then $\Tr_A X=\1_B$, so $\Tr X=b$ and $\Tr\eta_A=1$. Since $X>0$, $\eta_A>0$. Equation (2) gives $X\le b^2\eta_A\otimes\1_B$. Congruence by $\1_A\otimes\sigma_B^{1/2}$ yields (3). In particular $W_B=\sigma_B$. Crucially, $\eta_A$ need \emph{not} equal $\sigma_A$.
\shead{4. The score calculation}
If $r=1$, the mutual information vanishes identically. If $\delta=0$, the states coincide. Suppose $K>1$ and $0<\delta\le1-1/K$. Using (3), define
\[
 \widehat\sigma=\frac{KW-\sigma}{K-1},\qquad
 \tau=(1-\delta)\sigma+\delta\widehat\sigma=(1-c)\sigma+cW,\qquad c=\frac{K\delta}{K-1}.
\]
Both $\widehat\sigma$ and $\tau$ are density operators, $\tau_B=\sigma_B$, $0<c\le1$, and $\tau\ge(1-\delta)\sigma$. The two representations imply
\[
        cW\ \le\ \tau\ \le\bigl[(1-c)K+c\bigr]W=K(1-\delta)W .
\]
Define $G=-\log\tau+\log W$. Operator monotonicity of $\log$ gives
\begin{equation*}
 -\log[K(1-\delta)]\1\le G\le-\log c\,\1.\tag{4}
\end{equation*}
Here one adds the same Hermitian $\log W$ to the Loewner bounds on $-\log\tau$. No claim that $W$ and $\tau$ commute is needed.
\newpage
Write $\Delta=\rho-\sigma$ and $D(\alpha\|\beta)=\Tr\alpha(\log\alpha-\log\beta)$. The product logarithm is $\log W=(\log\eta_A)\otimes\1_B+\1_A\otimes\log\sigma_B$. The first term has pairing zero with $\Delta$ because $\Delta_A=0$. Expanding every relative entropy yields
\begin{align*}
 H(A|B)_\rho-H(A|B)_\sigma
 &=\Tr(\Delta G)+D(\sigma\|\tau)-D(\rho\|\tau)+D(\rho_B\|\sigma_B)\\
 &\le\Tr(\Delta G)+D(\sigma\|\tau).
\end{align*}
The inequality is data processing under $\Tr_A$, using $\tau_B=\sigma_B$. Both signs matter. Since $\tau\ge (1-\delta)\sigma$, operator monotonicity of $\log$ gives $D(\sigma\|\tau)\le-\log(1-\delta)$. The positive and negative parts of $\Delta$ each have trace $t=T(\rho,\sigma)$. Equation (4) gives
\[
 \Tr(\Delta G)\le t\log\frac{(K-1)(1-\delta)}{\delta}
                 \le \delta\log\frac{(K-1)(1-\delta)}{\delta};
\]
the coefficient is nonnegative throughout this branch. Adding the estimates proves
\[
 I(A:B)_\sigma-I(A:B)_\rho
 =H(A|B)_\rho-H(A|B)_\sigma
 \le h(\delta)+\delta\log(K-1).
\]
Switch $\rho$ and $\sigma$, constructing a fresh reference for the new anchor, for the reverse sign.

\textbf{Singular states.} Replace each state $\omega$ by
$(1-\epsilon)\omega+\epsilon\1_{AB}/(ab)$, $0<\epsilon<1$. These two regularizations remain equal on $A$, are positive definite (as are their marginals), and their trace distance is $(1-\epsilon)t\le\delta$. The common finite-dimensional entropy functions are continuous even at zero eigenvalues. Passing to the limit proves the stated estimate without ever applying a singular logarithm or comparing different marginal supports.

\textbf{Plateau.} Subadditivity and the Araki--Lieb inequality imply $0\le I(A:B)\le2\min(H(A),H(B))\le2\log r=\log K$. This proves the bound for $\delta\ge1-1/K$, including $\delta=1$.
\shead{5. Sharpness, including the strict unequal-marginal subclass}
Choose $r$-dimensional subspaces $A_0,B_0$, let $P$ be the identity on their product (zero elsewhere), and let $\Phi$ be a maximally entangled pure state projector there. Use
\[
  s=\min(\delta,1-1/K),\qquad
  \sigma=\Phi,\qquad \rho_s=(1-s)\Phi+\frac{s}{K-1}(P-\Phi).
\]
Both marginals coincide, and are maximally mixed on their respective subspaces. The trace distance is exactly $s$, and the nonzero eigenvalues of $\rho_s$ are $1-s$ and $s/(K-1)$ with multiplicity $K-1$ for $s>0$. Hence $I_\sigma-I_{\rho_s}=H(\rho_s)=g_K(\delta)$.
If $\rho_B\ne\sigma_B$ is required and $\delta>0$, take $s_n\uparrow s$ with $s_n<s$ and replace $\rho_{s_n}$ by $(1-\epsilon_n)\rho_{s_n}+\epsilon_n\pi_{A_0}\otimes|0\rangle\langle0|_{B_0}$, with $0<\epsilon_n<\delta-s_n$ tending to zero. Then $A$ agrees, $B$ differs, the distance is at most $s_n+\epsilon_n<\delta$, and continuity gives the same limiting gap.
\newpage
\shead{6. Exact channel-capacity corollary}
For channels $\N,\M:A\longrightarrow B$ let $d_\diamond(\N,\M)=\tfrac12\|\N-\M\|_\diamond$. If $d_\diamond(\N,\M)\le\delta$ then
\[
          |C_E(\N)-C_E(\M)|\le g_{r^2}(\delta).
          \tag{5}
\]
Indeed, for a common input $\alpha_A$ with purification $\psi_{RA}$, $\dim R=a$, the outputs on $RB$ of $\id_R\otimes\N$ and $\id_R\otimes\M$ agree on $R$ and have trace distance at most $\delta$. Use (1) for each input, then $|\sup f-\sup g|\le\sup|f-g|$ and the entanglement-assisted capacity formula [5].

\textbf{Equality in every dimension pair.} For $r\ge2$ put $s=\min(\delta,1-r^{-2})$ and use the $r$-dimensional identity channel $\id_r$ and depolarizing channel
\[
 \D_s(Z)=(1-p)Z+p(\Tr Z)\1_r/r,\qquad p=\frac{r^2}{r^2-1}s .
\]
As in [3, Rem.\ 3.2], $\D_s=(1-s)\id_r+s\mathcal R$, where $\mathcal R$ averages the $r^2-1$ nonidentity Weyl conjugations. The triangle inequality bounds the half diamond distance by $s$; the normalized maximally entangled input attains it because its nonidentity Weyl images are orthogonal to it. Unitary covariance and concavity of input mutual information show that the maximally mixed input optimizes both channels. Their normalized Choi states are exactly the two states in \S5, so their capacity difference is $g_{r^2}(\delta)$.

If $a\le b$, postcompose both square channels with the same isometric embedding into $B$, which preserves their distance and capacities (it has a CPTP left inverse). If $a>b$, let $V:\mathbb C^b\to\mathbb C^a$ be an isometry, $P_V=VV^*$, $E(Z)=VZV^*$, and
\[
 \C(X)=V^*XV+\Tr[(\1_a-P_V)X]\1_b/b .
\]
This is a channel (the second summand measures the complementary subspace and prepares the maximally mixed output), and $\C E=\id_b$. Precompose both square channels by $\C$. For every channel $L$ on $b$ levels, $C_E(L\C)\le C_E(L)$ by operational simulation, while $L=(L\C)E$ gives the reverse inequality. Contractivity under channel precomposition gives both inequalities for the diamond distance in exactly the same way. This extends the square extremizers to every $a>b$, including $b\nmid a$.
\shead{References}
{\small
[1] M. Berta, L. Lami, M. Tomamichel, \href{https://arxiv.org/html/2408.15226v2#S4}{\emph{Continuity of Entropies via Integral Representations}}, IEEE Trans. Inf. Theory (2025), \S IV immediately before (106).

[2] Y. Zhang, \href{https://doi.org/10.5281/zenodo.23253944}{\emph{A ternary counterexample to the proposed unrestricted mutual-information continuity bound}} (2026); \href{https://github.com/mxym/math/blob/main/preprints/mutual-information-continuity-2026-10/README.md}{project scope}. Both marginals change.

[3] M. Berta, P. Costa Rico, G. Kossmann, L. Lami, J. A. Zeiss, \href{https://arxiv.org/html/2607.24687v2}{\emph{Sharp continuity of quantum conditional entropy}}, arXiv:2607.24687v2 (28 Sep 2026), Thm.\ 1.1, Prop.\ 2.1, Prop.\ 3.1 and Rem.\ 3.2.

[4] M. Berta, M. Christandl, R. Renner, \href{https://arxiv.org/html/0912.3805v3#A2.SS3}{\emph{The Quantum Reverse Shannon Theorem Based on One-Shot Information Theory}}, Commun. Math. Phys. 306 (2011), Lem.\ B.7.

[5] C. H. Bennett, P. W. Shor, J. A. Smolin, A. V. Thapliyal, \href{https://arxiv.org/abs/quant-ph/0106052}{\emph{Entanglement-assisted capacity of a quantum channel and the reverse Shannon theorem}}, IEEE Trans. Inf. Theory 48 (2002).

[6] M. E. Shirokov, \href{https://arxiv.org/abs/1512.09047}{\emph{Tight uniform continuity bounds for the quantum conditional mutual information, for the Holevo quantity, and for capacities of quantum channels}}, J. Math. Phys. 58, 102202 (2017), Cor.\ 1 and Prop.\ 10.}
\end{document}
